\documentclass[12pt]{article}
\usepackage{../../resource/shortex}
\usepackage{mathtools,amssymb}
\usepackage{enumitem}
\usepackage{datetime}
\usepackage{fancyhdr}
\usepackage{ifthen}
\usepackage{pifont}
\usepackage{mathrsfs}
\usepackage[normalem]{ulem}
%\RequirePackage[usenames,dvipsnames]{color}


%% showing solutions 
\newboolean{showsols}
\setboolean{showsols}{false}
\newcommand{\showsolutions}{\setboolean{showsols}{true}}
\newlength{\solspace}
\setlength{\solspace}{16em}
\newcommand{\nosolspace}{\setlength{\solspace}{0em}}

\newcommand{\solution}[2]{\ifthenelse{\boolean{showsols}}{{\par\textcolor{red}{Rubric criterion: #1}\newline\textcolor{blue}{#2}}}{\vspace{\solspace}}}
\newcommand{\npforprint}{\ifthenelse{\boolean{showsols}}{}{\newpage}}
\newcommand{\vsforprint}[1]{\ifthenelse{\boolean{showsols}}{}{\vspace{#1}}}

%% grading
\newcommand{\grades}[1]{~\newline\noindent\framebox[\textwidth][c]{E\quad \quad \quad \quad S \quad \quad \quad \quad R \quad \quad \quad \quad N}}

\newcommand{\questiongrades}{\noindent\emph{Each item is marked $\checkmark$, $\checkmark\!-$, or \ding{55}:}\begin{center}\emph{$\checkmark$ = capable  $~~$ $\checkmark\!-$ = mostly capable  $~~$ \ding{55} = not yet}\end{center} }

\newcommand{\triAssessments}{$\checkmark\quad\checkmark\!-\quad$\ding{55}\quad}
\newcommand{\biAssessments}{$\checkmark\quad$\ding{55}\quad}

%% line for writing name and BUID
\newcommand{\nameBUID}{\noindent\textbf{Name:}\underline{\phantom{XXXXXXXXXXXXXXXXXXX}}\hfill\textbf{BUID:}\underline{\phantom{XXXXXXXXXXXXX}}\newline}


%% formatting 
%\allowdisplaybreaks[3]
\renewcommand{\headrulewidth}{.4mm} % header line width
\renewcommand{\arraystretch}{1.25}


%% semester label, read from the "Semester:" line of important_dates.md
%% (path is relative to a document in materials/ModuleX/<folder>/)
\ExplSyntaxOn
\tl_new:N \g_teaching_semester_tl
\ior_new:N \g_teaching_dates_ior
\file_if_exist:nTF { ../../instructor_inputs/important_dates.md }
  {
    \ior_open:Nn \g_teaching_dates_ior { ../../instructor_inputs/important_dates.md }
    \ior_str_map_inline:Nn \g_teaching_dates_ior
      {
        \regex_extract_once:nnNT { \A \s* Semester: \s* (.*?) \s* \Z } {#1} \l_tmpa_seq
          {
            \tl_gset:Nx \g_teaching_semester_tl { \seq_item:Nn \l_tmpa_seq { 2 } }
            \ior_map_break:
          }
      }
    \ior_close:N \g_teaching_dates_ior
  }
  {
    \msg_new:nnn { teaching } { no-dates }
      { ../../instructor_inputs/important_dates.md~not~found;~semester~header~left~blank. }
    \msg_warning:nn { teaching } { no-dates }
  }
\cs_new:Npn \semester { \tl_use:N \g_teaching_semester_tl }
\ExplSyntaxOff

%% page style
\pagestyle{fancy}
\fancyhf{}
\rhead{\semester}
\lhead{CAS MA 213: Basic Statistics and Probability}
%\rfoot{\thepage}

\definecolor{checkincolor}{RGB}{232,222,250}
\newcommand{\calloutbox}[3]{%
  \vspace{4pt}\noindent
  \colorbox{#1}{\parbox{0.96\linewidth}{\textbf{#2}}}\\*[2pt]
  \noindent\fbox{\begin{minipage}{0.94\linewidth}\vspace{2pt}#3\vspace{2pt}\end{minipage}}
  \vspace{6pt}
}
\newcommand{\checkinbox}[1]{%
  \calloutbox{checkincolor}{CHECK IN: Call over an instructor}{%
    \begin{minipage}[c]{0.72\linewidth}#1\end{minipage}\hfill
    \begin{minipage}[c]{0.2\linewidth}\centering
      \fbox{\rule{0pt}{0.45in}\hspace{0.9in}}\\[2pt]
      {\small Instructor initials}
    \end{minipage}}}

%--------------------
\begin{document}

\begin{center}
\textbf{\Large Project 1, Part 2 \\ Lab Checklist}
\end{center}

\nameBUID

\begin{center}
\begin{tabular}{ll}
\toprule
\textbf{Item} & \textbf{Date} \\
\midrule
Lab P1\_2 (today) & Fri, Oct 16 \\
\textbf{Deliverable 3 (recorded presentation) due} & \textbf{Thu, Oct 22, 9:00 PM} \\
\bottomrule
\end{tabular}
\end{center}

\subsection*{Goal for Today}
Turn your Deliverable 2 slides and R code into a finished recorded presentation. By the end of lab, you should have checked in with an instructor about your group's R code, your group should know what will be said on each slide and who will say it, and you should have a plan for recording.

\vspace{6pt}
\noindent If you don't get to step 6 with 10 minutes left, stop where you are and move straight to step 6.

\begin{enumerate}
    \item[$\square$] \textbf{1. Review your Deliverable 2 feedback.} Read it together, list every change it asks for, and give each change one owner.

    \item[$\square$] \textbf{2. Instructor R code check-in.} Keep working on the other steps until the instructor gets to you.

    \checkinbox{R Check-in.}

    \pagebreak
    \item[$\square$] \textbf{3. Revise the R code.}
    \begin{itemize}[nosep]
        \item[$\square$] Changes from the feedback (and from the check-in) are made
        \item[$\square$] Script runs from top to bottom without errors
        \item[$\square$] Data file is included in the submission
        \item[$\square$] Uses only functions covered in class, tutorials, or labs, or in the standard R distribution
        \item[$\square$] Appropriate comments are included to explain what each section of code does
    \end{itemize}

    \item[$\square$] \textbf{4. Revise the slides.}
    \begin{itemize}[nosep]
        \item[$\square$] Changes from the feedback are made
        \item[$\square$] Numerical: histogram, boxplot, scatter plot, and summary statistics table, each clearly labeled
        \item[$\square$] Categorical: bar plot and a frequency or contingency table, each clearly labeled
        \item[$\square$] Data description mentions missing data and/or outliers
    \end{itemize}

    \item[$\square$] \textbf{5. Plan the presentation.} For each slide:
    \begin{itemize}[nosep]
        \item[$\square$] Decide what you will say: the main point, plus what the audience should notice in the plot or table (trend or association, and the shape, center, spread, and outliers of a distribution)
        \item[$\square$] Decide who will say it
        \item[$\square$] Consider removing extra text from the slide if you are going to say it
    \end{itemize}
    Aim for about 5 minutes total --- roughly 30--60 seconds per slide.

    \item[$\square$] \textbf{6. Plan the recording and the Q\&A.}
    \begin{itemize}[nosep]
        \item[$\square$] Picked a date and time to record (before the Deliverable 3 deadline)
        \item[$\square$] Decided who records, who edits the video, and who submits it and posts it to Blackboard
    \end{itemize}
\end{enumerate}

\end{document}
